\section{Introduction}

CLIP's zero-shot accuracy is strong on natural-image classification but collapses on
out-of-di\-stri\-bution domains: the benchmark that motivates this project reports 23.7\% on
EuroSAT and 15.7\% on Aircraft~\cite{sheng2025illusion}. Our own zero-shot measurements
(Table~\ref{tab:main}) are higher -- 42.06\% and 23.85\% -- because we evaluate the CoOp test
splits with a \texttt{ViT-B/16} backbone rather than that paper's exact setup; the collapse
relative to in-distribution accuracy is the point, and it reproduces. Test-time adaptation (TTA) methods
close much of this gap without any labels or backpropagation. TDA~\cite{karmanov2024tda}, in
particular, keeps a small cache of features from test images the model is confident about and
reuses that cache to refine later predictions -- a forward-pass-only method with no training
loop.

Accuracy, however, is not the only axis that matters for a system meant to be trusted. A
model that is right 80\% of the time but reports 95\% confidence on every prediction is
\emph{worse} to deploy than one that is honest about its 80\%, because a downstream user or
pipeline has no way to tell a confident-and-correct prediction from a confident-and-wrong one.
This is exactly the failure mode \citet{sheng2025illusion} document in TDA: its cache rewards
confident predictions with more confidence regardless of correctness, so accuracy gains arrive
bundled with a calibration loss (average ECE rising 5.70\% $\rightarrow$ 9.21\% across their
benchmark).

This project asks whether that calibration gap can be closed without retraining and without
labels, by changing three seams of TDA's cache loop:

\begin{enumerate}
    \item \textbf{Margin-based cache admission.} Admit a test sample to the cache by the gap
        between its top-1 and top-2 class scores, instead of by prediction entropy, so the
        cache does not accumulate merely \emph{loud} predictions.
    \item \textbf{Probabilistic fusion.} Combine CLIP's and the cache's evidence as a convex
        mixture of probability distributions, instead of adding raw logits, so agreement
        between the two sources cannot manufacture confidence that neither source had on its
        own.
    \item \textbf{Leave-one-out (LOO) temperature scaling.} Estimate the cache's own accuracy
        without ground truth, via leave-one-out nearest-neighbour agreement among cached
        items, and pick a softmax temperature that makes the model's mean confidence match
        that label-free estimate.
    \end{enumerate}

Each change is independently switchable from configuration, which is what makes the ablation
in Section~\ref{sec:results} possible -- and the ablation is where this project's real finding
lives. \textbf{The three changes together do not improve calibration: they make it
dramatically worse.} Run in isolation, however, one of the three delivers exactly the
calibration improvement this project set out to find, at no accuracy cost. We report both
results without softening either one.
